% Generated draft: requires mathematical and source review.
\documentclass[11pt]{article}
\usepackage[margin=25mm]{geometry}
\usepackage{fontspec}
\setmainfont{Latin Modern Roman}
\usepackage{amsmath,amssymb,mathtools,unicode-math}
\setmathfont{Latin Modern Math}
\usepackage{enumitem,needspace,xurl,hyperref}
\setlength{\parindent}{0pt}
\setlength{\parskip}{5pt}
\setlength{\emergencystretch}{3em}
\newcommand{\N}{\mathbb{N}}
\newcommand{\Z}{\mathbb{Z}}
\newcommand{\Q}{\mathbb{Q}}
\newcommand{\R}{\mathbb{R}}
\newcommand{\C}{\mathbb{C}}
\newcommand{\F}{\mathbb{F}}
\newcommand{\A}{\mathbb{A}}
\newcommand{\PP}{\mathbb{P}}
\newcommand{\E}{\mathbb{E}}
\newcommand{\eps}{\varepsilon}
\newcommand{\dd}{\,\mathrm{d}}
\begin{document}
\section*{UnsolvedMath ID: 30000705}
\textbf{Source problem code:} catalogue entry
\subsection{Definitions and mathematical statement}
Let $G$ be a semigroup under composition of rational self-maps of the Riemann sphere $\widehat\C$, using the usual degree-at-least-two convention. Its Fatou set is the largest open set on which the family $G$ is normal, and $J(G)$ is its complement. For $h\in G$, let $J_h$ be the Julia set of the single map, defined using its iterates. A set $E$ is completely invariant under $G$ if $h(E)=E=h^{-1}(E)$ for every $h\in G$. Let $E(G)$ be the source's smallest closed completely invariant subset containing at least three points.

The first conjecture is
\[
\left[\exists f,g\in G:\ J_f\ne J_g\right]
\ \text{and}\ E(G)\ne\widehat\C
\quad\Longrightarrow\quad
E(G)\text{ is Möbius-equivalent to a line segment or a circle}.
\]
Möbius equivalence means that some fractional-linear automorphism of $\widehat\C$ maps $E(G)$ onto the indicated set. The nonspherical hypothesis concerns $E(G)$, not $J(G)$.

The second conjecture has its own hypotheses: assume every map in $G$ is a polynomial, assume $J_h=J(G)$ for some $h\in G$, and assume $J(G)$ is not a line segment. It asks whether
\[
\forall f,g\in G,\qquad J_f=J_g,
\]
and, as the source's asserted further consequence, whether $G$ is nearly abelian. The different-Julia-set hypothesis of the first question is not carried over to this second question.

Near abelianness is ordinarily expressed through a controlled family $\Phi$ of Möbius transformations with relations
\[
\forall f,g\in G\ \exists\phi\in\Phi,\qquad
f\circ g=\phi\circ g\circ f,
\]
with additional compactness and Fatou-set preservation conditions in the intended definition. The extract gives no precise definition of $\Phi$ or near abelianness; it is not replaced by ordinary commutativity or by the displayed relation alone. It also does not explicitly state degree conventions, finite generation, or the construction guaranteeing the stated minimal invariant set. No finite-generation or extra dynamical hypothesis is added.
\subsection{Short English statement}
First classify proper completely invariant sets when two rational maps in the semigroup have different Julia sets. Separately, ask whether a polynomial semigroup whose Julia set equals that of one element, and is not a segment, must have equal elementwise Julia sets and be nearly abelian.
\subsection{Sources}
\begin{itemize}
\item Supplied report: \emph{Normal Families and Complex Dynamics} (2007), Oberwolfach Reports, DOI 10.4171/owr/2007/09, pp. 487--548; proposed year 2007. \url{https://doi.org/10.4171/owr/2007/09}.
\item Supplied reference labelled ``counterexamples and remaining conjecture'': \url{https://arxiv.org/abs/0708.3434}.
\item Local UnsolvedMath registry, permanent ID 30000705, catalogue entry; created 25 August 2026, updated 16 September 2026. Its corrected\_verified note describes editorial cleanup. The supplied 21 August 2026 literature triage labels the record open and says later work retains the non-spherical line-segment-or-circle conjecture as unresolved. That assessment addresses the first classification question explicitly and does not independently establish the status of the second. No updated literature assessment is asserted here.
\end{itemize}
\end{document}
